\documentclass[11pt]{article}
% Shared preamble for the hanzo-kernel Paper 07 series.
% One style, one vocabulary, inputted by every paper so the series reads as one voice.
\usepackage[margin=1in]{geometry}
\usepackage{booktabs}
\usepackage{amsmath}
\usepackage{amssymb}
\usepackage{graphicx}
\usepackage{xcolor}
\usepackage{hyperref}
\hypersetup{colorlinks=true, linkcolor=blue, urlcolor=blue, citecolor=blue}
\usepackage{enumitem}
\usepackage{listings}
\usepackage{array}

\definecolor{kw}{rgb}{0.13,0.29,0.53}
\definecolor{cmt}{rgb}{0.40,0.40,0.40}
\definecolor{str}{rgb}{0.16,0.50,0.16}
\lstset{
  basicstyle=\ttfamily\footnotesize,
  keywordstyle=\color{kw}\bfseries,
  commentstyle=\color{cmt}\itshape,
  stringstyle=\color{str},
  showstringspaces=false,
  breaklines=true,
  columns=fullflexible,
  frame=single,
  framesep=4pt,
}
\lstdefinelanguage{Rust}{
  morekeywords={fn,let,mut,pub,struct,impl,trait,enum,match,if,else,for,while,loop,return,use,mod,crate,self,super,async,await,where,type,const,static,unsafe,ref,move,dyn,as,true,false},
  morecomment=[l]{//}, morecomment=[s]{/*}{*/}, morestring=[b]", sensitive=true,
}

\newcommand{\x}{\ensuremath{\times}}
\newcommand{\dpa}{\texttt{dp4a}}
\newcommand{\hk}{\texttt{hanzo-kernel}}
\newcommand{\code}[1]{\texttt{#1}}

% Absorb minor prose overfulls without rivers; keep line-breaking sane.
\tolerance=800
\emergencystretch=3em


\usepackage{amsthm}
\newtheorem{proposition}{Proposition}
\newtheorem{lemma}{Lemma}
\newtheorem{definition}{Definition}
\newtheorem{remark}{Remark}

\newcommand{\E}{\mathbb{E}}
\newcommand{\R}{\mathbb{R}}
\newcommand{\avg}[1]{\langle #1 \rangle}
\newcommand{\heff}{h_{\mathrm{eff}}}
\DeclareMathOperator{\artanh}{artanh}
\DeclareMathOperator{\sech}{sech}
\DeclareMathOperator{\sign}{sign}

\title{\textbf{Critical Alignment of Agentic Fleets:\\ Mean-Field Games and Ferromagnetic Order}}
\author{Hanzo AI --- ML Systems}
\date{hanzo-agentic Paper 01 --- the \emph{Continually Improving Fleets} series}

\begin{document}
\maketitle

\begin{abstract}
A fleet of agents pointed at one problem has two failure modes and no natural middle. Left independent it
\emph{fragments}: incoherent, duplicated, wasted work. Coupled too hard it \emph{collapses} into groupthink:
a single consensus, possibly wrong, that no member can break. We want the middle --- alignment that is
coherent \emph{and} correctable --- to \emph{emerge} from a control law rather than be commanded. We supply
that law by an exact structural correspondence. A mean-field game (MFG) closes the Hamilton--Jacobi--Bellman
best-response against the Fokker--Planck population-flow on a self-consistent fixed point
$(u^{*},m^{*})$: Nash against the distribution. Curie--Weiss ferromagnetism is that \emph{same} loop reduced
to a scalar order parameter $m$: each spin best-responds to the mean field $\heff = Jm + h$, and consistency
gives $m=\tanh(\beta(Jm+h))$. We prove Curie--Weiss is a static potential MFG on $\{\pm1\}$ with quadratic
interaction (so $\mathrm{MFG}\supseteq$ Curie--Weiss), that a nonzero $m$ bifurcates at $T_c=J$
($\beta J>1$), and that the susceptibility $\chi=\beta/(1-\beta J)$ diverges as $T\to T_c^{+}$ (Curie--Weiss
law): near criticality an infinitesimal field produces macroscopic aligned response. This hands us four
knobs --- order parameter (fleet alignment), coupling (shared context), field (the directive), temperature
(independent-reasoning diversity) --- one prescription (operate just below $T_c$; align by the field, not the
coupling), and two named failure modes (groupthink, mob). We then ground it twice in shipped Hanzo systems:
GRPO's group-relative advantage \emph{is} the mean field, which turns the criticality prescription into a
training-hyperparameter schedule; and the agentic harness's adversarial critic serves triple duty ---
ferromagnetic fluctuation, RL reward, and selection signal --- closing a model$\leftrightarrow$fleet
co-evolution flywheel. We keep the accounting honest: this is a control lens, not a literal Hamiltonian; the
physics buys the right variables and the phase-transition intuition, not free rigor.
\end{abstract}

\section{The pain: fragment or stampede, nothing in between}
\label{sec:pain}

Run $N$ agents at one hard problem and dial the single knob you obviously have --- how much they listen to each
other --- and you traverse a line with a bad end at each extreme and, apparently, nothing usable in the middle.
Turn the coupling to zero and the fleet \emph{fragments}: $N$ lonely searches that rediscover each other's dead
ends, no one's partial win reaching anyone else, throughput linear in $N$ and insight sublinear. Turn the
coupling up and the fleet \emph{stampedes}: it locks onto a running consensus and amplifies it, and if that
consensus is wrong there is no longer any member far enough outside it to say so. The second failure is the
worse one because it \emph{looks} like success --- full agreement, high confidence, clean logs --- right up to
the moment the shipped artifact fails in production. This is groupthink, and once a fleet is in it a single
correct dissent cannot escape: the basin is too deep for one voice to climb out of.

What we actually want is the state a good research group occupies when it is working: enough shared context that
one person's insight propagates and the group converges on \emph{a} direction, but enough surviving
independence that a single member who spots the flaw can still turn the whole group. Call it
\emph{coherent-but-correctable}. The thesis of this paper is that this state is not a fragile sweet spot you
tune by hand and hope to hold; it is a \emph{critical point}, it has a name in physics, and the same
mathematics that names it also tells you which knob buys coherence, which knob buys correctness, and where to
sit so that both survive. We make the correspondence exact, extract the control law, and then show the law is
already running --- unnamed --- inside Hanzo's model-training loop and its agent harness. This is the inner
alignment layer of the \emph{Continually Improving Fleets} stack whose outer evolutionary loop is the companion
\emph{Evolutionary Orchestration} \cite{hanzo-evolutionary-orchestration}.

\section{The isomorphism: a mean-field game with a scalar order parameter}
\label{sec:iso}

\subsection{The two loops are one loop}

A mean-field game \cite{lasrylions2007} models a continuum of identical rational agents, each choosing a
control to minimise its own cost, where that cost depends on the \emph{distribution} $m$ of everyone else. Its
equilibrium is the coupled pair
\begin{align}
-\partial_t u - \nu\,\Delta u + H\!\left(x,\nabla u\right) &= F(x,m), &&\text{(HJB: best response)} \label{eq:hjb}\\
\partial_t m - \nu\,\Delta m - \operatorname{div}\!\left(m\,\nabla_p H(x,\nabla u)\right) &= 0, &&\text{(FPK: population flow)} \label{eq:fpk}
\end{align}
closed at a fixed point $(u^{*},m^{*})$: the value function $u^{*}$ is each agent's optimal response to the
population $m^{*}$ (Hamilton--Jacobi--Bellman, run backward), and the feedback control read off from $u^{*}$
transports the population so that its own law reproduces $m^{*}$ (Fokker--Planck, run forward). No agent can do
better by deviating from the control that is optimal against $m^{*}$; the equilibrium is \emph{Nash against the
distribution}. The structural content is a single self-consistency: best-response $\otimes$
population-consistency, closed.

Curie--Weiss ferromagnetism is that self-consistency with the population measure collapsed to one number. Take
$N$ spins $\sigma_i\in\{\pm1\}$ with the all-to-all Curie--Weiss energy
\begin{equation}
\mathcal{E}(\sigma) \;=\; -\frac{J}{2N}\sum_{i,j}\sigma_i\sigma_j \;-\; h\sum_i \sigma_i ,
\label{eq:cwenergy}
\end{equation}
so that each spin feels the same \emph{mean field} $\heff = Jm+h$ generated by the average magnetisation
$m=\frac1N\sum_i\sigma_i$ and an external field $h$ \cite{baxter1982}. Because every spin couples only through
the scalar $m$, the measure-valued object $m^{*}$ of the MFG reduces to a scalar order parameter, and the two
PDEs of \eqref{eq:hjb}--\eqref{eq:fpk} reduce to two lines.

\begin{definition}[Curie--Weiss self-consistency]
\label{def:cw}
The \emph{best response} of one spin to a field $\heff$ is the single-site Gibbs law
$p_\beta(\sigma)\propto e^{\beta\heff\sigma}$ at inverse temperature $\beta=1/T$, whose mean is
$\avg{\sigma}=\tanh(\beta\heff)$. \emph{Population consistency} demands the field-generating average equal that
best response, $m=\avg{\sigma}$. The pair closes at the fixed point
\begin{equation}
m^{*} \;=\; \tanh\!\bigl(\beta(Jm^{*}+h)\bigr).
\label{eq:selfconsistent}
\end{equation}
\end{definition}

\begin{proposition}[Ferromagnetism is a scalar-order-parameter potential MFG]
\label{prop:isomorphism}
The Curie--Weiss self-consistency \eqref{eq:selfconsistent} is exactly the fixed-point condition of a mean-field
game whose state space is $\{\pm1\}$, whose per-agent best response is the free-energy--minimising single-site
Gibbs law against the mean field $\heff=Jm+h$, and whose population-consistency map is $m=\avg{\sigma}$.
The game is a \emph{potential} game: there is one common potential, the mean-field free energy per spin
\begin{equation}
f(m) \;=\; -\tfrac12 J m^2 \;-\; h m \;-\; T\,\mathcal{H}(m), \qquad
\mathcal{H}(m) = -\tfrac{1+m}{2}\ln\tfrac{1+m}{2} - \tfrac{1-m}{2}\ln\tfrac{1-m}{2},
\label{eq:freeenergy}
\end{equation}
whose stationary points are precisely the self-consistent $m^{*}$. Hence Curie--Weiss ferromagnetism is the
static, two-state, quadratic-interaction special case of a Lasry--Lions MFG: $\mathrm{MFG}\supseteq$
Curie--Weiss.
\end{proposition}

\begin{proof}
By the Gibbs variational principle, the single-site law minimising the site free energy
$\E_p[-\heff\sigma]-T\,\mathcal{H}[p]$ over distributions on $\{\pm1\}$ is the Boltzmann law
$p_\beta(\sigma)\propto e^{\beta\heff\sigma}$, unique because the objective is strictly convex in $p$; its mean
is $\tanh(\beta\heff)$. This is the reduction of the HJB best-response \eqref{eq:hjb}: with a static single-step
decision and no spatial transport, the value function collapses to the site free energy and its minimiser is the
Gibbs law. Imposing the population-consistency $m=\avg{\sigma}$ --- the reduction of the FPK stationarity
\eqref{eq:fpk}, since the interaction enters only through the mean $m$ so the stationary population is pinned by
its own average --- yields \eqref{eq:selfconsistent}.

For the potential structure, differentiate \eqref{eq:freeenergy} using $\mathcal{H}'(m)=-\artanh(m)$:
\[
f'(m) \;=\; -Jm - h + T\,\artanh(m).
\]
Setting $f'(m)=0$ gives $\artanh(m)=\beta(Jm+h)$, i.e. $m=\tanh(\beta(Jm+h))$, which is
\eqref{eq:selfconsistent}. The coupling term $-Jm\sigma$ of \eqref{eq:cwenergy} derives from the shared
quadratic $-\tfrac12 Jm^2$: every agent's incentive depends on the others only through the common scalar $m$,
which is exactly the defining property of a potential game --- one function whose descent all agents follow.
Thus the MFG is potential with potential $f$, and its equilibria are the critical points of $f$.
\end{proof}

The reading is: the ferromagnet is not \emph{like} a mean-field game; it \emph{is} one, with the whole
population distribution compressed to a single order parameter because the agents interact only through their
average. Everything below is therefore a statement about MFG equilibria that we are entitled to compute in one
scalar variable.

\subsection{Where order switches on: the critical temperature}

\begin{proposition}[Bifurcation at $T_c=J$]
\label{prop:tc}
Let $h=0$. Then $m=0$ solves \eqref{eq:selfconsistent} for every $T$, and a nonzero solution $m^{*}\neq0$
exists if and only if $\beta J>1$, i.e. $T<T_c$ with the critical temperature $T_c=J$. Near the transition the
spontaneous magnetisation grows as $m^{*}\sim\sqrt{3(\beta J-1)}\,/(\beta J)^{3/2}$, the mean-field square-root
law.
\end{proposition}

\begin{proof}
At $h=0$, \eqref{eq:selfconsistent} reads $m=\tanh(\beta J m)$; $m=0$ is always a root. Expand
$\tanh(x)=x-\tfrac13 x^3+O(x^5)$ with $x=\beta J m$:
\[
m \;=\; \beta J m \;-\; \tfrac13(\beta J m)^3 + O(m^5)
\quad\Longrightarrow\quad
(1-\beta J)\,m \;=\; -\tfrac13(\beta J)^3 m^3 + O(m^5).
\]
If $\beta J<1$ the left coefficient is positive and the right side is nonpositive for small $m$, forcing $m=0$:
the only root near the origin is trivial. If $\beta J>1$ the left coefficient is negative and a nonzero balance
appears, $m^{*2}=3(\beta J-1)/(\beta J)^3+O(\beta J-1)$, giving the stated $m^{*}$. The exchange of stability
happens at $\beta J=1$, i.e. $T=J$; equivalently the slope of the return map $\tfrac{d}{dm}\tanh(\beta Jm)|_{0}
=\beta J$ crosses unity there. Hence $T_c=J$.
\end{proof}

Below $T_c$ the fleet spontaneously orders; above it, it does not. The transition is sharp in $m$ but --- as the
next proposition shows --- announced from \emph{above} by a diverging responsiveness.

\subsection{Why criticality is the useful place: divergent susceptibility}

\begin{proposition}[Curie--Weiss law]
\label{prop:chi}
In the disordered phase ($T>T_c$, so $m^{*}=0$ at $h=0$) the static susceptibility is
\begin{equation}
\chi \;=\; \left.\frac{\partial m}{\partial h}\right|_{h=0} \;=\; \frac{\beta}{1-\beta J} \;=\; \frac{1}{T-T_c},
\label{eq:chi}
\end{equation}
which diverges as $T\to T_c^{+}$. An infinitesimal field therefore produces a macroscopic aligned response near
criticality. Deep in the ordered phase the response instead \emph{vanishes}: with $m^{*}\to1$ as $T\to0$, the
susceptibility $\chi=\beta(1-m^{*2})/\bigl(1-\beta J(1-m^{*2})\bigr)\to0$.
\end{proposition}

\begin{proof}
Differentiate \eqref{eq:selfconsistent} in $h$, writing $s=\sech^2\!\bigl(\beta(Jm^{*}+h)\bigr)$:
\[
\chi \;=\; s\,\beta\,(J\chi+1)
\quad\Longrightarrow\quad
\chi \;=\; \frac{s\,\beta}{1 - s\,\beta J}.
\]
At $h=0$, $T>T_c$ we have $m^{*}=0$ and $s=\sech^2(0)=1$, giving $\chi=\beta/(1-\beta J)$; substituting
$\beta=1/T$ and $J=T_c$ yields $\chi=1/(T-T_c)$, which $\to+\infty$ as $T\downarrow T_c$. In the ordered phase
$s=\sech^2(\beta Jm^{*})=1-\tanh^2(\beta Jm^{*})=1-m^{*2}$ by \eqref{eq:selfconsistent}; as $T\to0$,
$m^{*}\to1$, so $s\to0$ and $\chi\to0$.
\end{proof}

Proposition~\ref{prop:chi} is the crux of the whole control story. Just above $T_c$ the fleet is still
disordered ($m\approx0$) yet \emph{maximally suggestible}: a vanishingly small correct nudge swings a
macroscopic fraction of it into alignment. Just below $T_c$ it is coherent yet still soft --- $\chi$ finite but
large, fluctuations alive. Deep below $T_c$ it is rigid --- $\chi\to0$, no nudge moves it, no critic can flip
it. Coherent-but-correctable is not a hand-tuned compromise; it is the neighbourhood of the critical point,
where susceptibility is large in both senses: the fleet can be aimed cheaply and un-stuck cheaply.

\section{The control law: four knobs, one operating point}
\label{sec:control}

Proposition~\ref{prop:isomorphism} promoted a physics model to a control law by naming its variables. Read as a
fleet controller, the four parameters are:

\begin{center}
\begin{tabular}{@{}lll@{}}
\toprule
\textbf{Physics} & \textbf{Fleet} & \textbf{Realisation} \\
\midrule
order parameter $m$ & fleet alignment & how concentrated the agents' behaviour/answer is \\
coupling $J$ & shared context & canonical recipe, shared memory, running consensus \\
field $h$ & the directive & the objective $+$ invariants the fleet must honour \\
temperature $T$ & independent reasoning & adversarial ``try to refute'', distinct lenses, isolation \\
\bottomrule
\end{tabular}
\end{center}

The coupling $J$ is everything that makes one agent heed the others: a shared canonical recipe, a shared memory,
a running consensus maintained across the fleet --- in Hanzo, the distributed consensus-and-compute substrate
that carries that shared state \cite{hanzo-consensus-zap-compute}. The temperature $T$ is everything that keeps
an agent thinking for itself: an explicit adversarial instruction to refute the current consensus, a distinct
analytical lens, or plain isolation from the others' partial answers. The field $h$ is the directive plus the
non-negotiable invariants. The two failure modes of Section~\ref{sec:pain} are now precise phases:

\begin{itemize}[leftmargin=1.4em,itemsep=2pt]
\item \textbf{Groupthink} is the deep ordered phase, $T\ll T_c$, $m\to1$. The adversarial value of independent
perspectives has collapsed ($\chi\to0$, Proposition~\ref{prop:chi}); a wrong consensus is a global minimum of
$f$ that no single dissent can climb out of. Agreement is total and correction is impossible.
\item \textbf{Mob} is the disordered phase, $T\gg T_c$, $m\approx0$. There is no shared direction at all; the
fleet is $N$ uncorrelated walks (Section~\ref{sec:pain}'s fragmentation).
\end{itemize}

The prescription is to operate \emph{just below} $T_c$ --- criticality, the edge of chaos --- where the
correlation length and susceptibility are largest: the fleet coheres, but fluctuations survive so a single
critic can still break a bad basin. And there is a sharp division of labour between the two knobs that produce
order.

\begin{proposition}[Coherence comes from coupling; aim comes from the field]
\label{prop:aim}
Fix $\beta J>1$ ($T<T_c$). At $h=0$ the potential \eqref{eq:freeenergy} is even, $f(m)=f(-m)$, so it has two
degenerate global minima $\pm m^{*}$ with $m^{*}>0$: the ordered fleet is coherent ($|m^{*}|>0$) but its
direction is undetermined, chosen by spontaneous symmetry breaking with probability $\tfrac12$ each. For
$h\neq0$ the degeneracy is lifted: the global minimiser satisfies $\sign(m^{*})=\sign(h)$, so the field alone
selects which direction the fleet agrees on. Near criticality the selection is arbitrarily cheap: by
Proposition~\ref{prop:chi}, $\chi=1/(T-T_c)\to\infty$ as $T\to T_c^{+}$, so an infinitesimal correct field
yields a macroscopic correct alignment.
\end{proposition}

\begin{proof}
At $h=0$ the term $-hm$ vanishes and both $-\tfrac12Jm^2$ and $\mathcal{H}(m)$ are even, so $f$ is even;
Proposition~\ref{prop:tc} gives two nonzero minima $\pm m^{*}$, equal by evenness. For $h\neq0$ the added linear
term $-hm$ strictly lowers $f$ at the aligned sign and raises it at the opposite sign, so
$f(\sign(h)\,|m^{*}|)<f(-\sign(h)\,|m^{*}|)$ and the global minimiser takes $\sign(m^{*})=\sign(h)$. The
cheapness of the selection near $T_c$ is Proposition~\ref{prop:chi}.
\end{proof}

The lesson is operational: \textbf{align by the field, not the coupling.} Cranking the coupling $J$ raises
$T_c$ and grows $|m^{*}|$ --- it buys \emph{coherence}, but it does not say \emph{what} the fleet should cohere
\emph{on}; at $h=0$ a strongly coupled fleet is exactly as likely to lock onto the wrong answer as the right
one, and then it is stuck there (that is groupthink, manufactured by coupling). Only the field $h$ --- the
directive and its invariants --- selects the direction, and at criticality a small, correct field is enough to
aim a coherent fleet at the objective. A controller that tries to force agreement by raising coupling is
building a rigid consensus with no guarantee it points anywhere good; a controller that holds the fleet at
criticality and aims it with a clear directive gets coherence \emph{and} correctness with cheap, correctable
alignment.

\section{Grounding I: the GRPO group is a Curie--Weiss ensemble}
\label{sec:grpo}

The control law is not an analogy imported from outside the stack; the mean field is already live in Hanzo's
trainer. Group Relative Policy Optimization (GRPO) \cite{shao2024grpo}, implemented and tested in the native
Rust training stack \cite{hanzo-native-training} (\code{ml/hanzo-training/src/grpo/}), samples for each prompt a
\emph{group} of $G$ rollouts, scores them, and forms the group-relative advantage
\begin{equation}
A_i \;=\; \frac{r_i - \operatorname{mean}(r)}{\operatorname{std}(r)}, \qquad
\operatorname{mean}(r)=\frac1G\sum_{j=1}^G r_j,
\label{eq:grpo}
\end{equation}
before a length-normalised policy-gradient step with an optional low-variance $k3$ KL penalty to a frozen
reference, $k3=\exp\Delta-\Delta-1$ with $\Delta=\log(\pi_{\mathrm{ref}}/\pi_\theta)$
(\code{math::batch\_advantages}, unbiased std with $\mathrm{ddof}{=}1$; tested). Equation \eqref{eq:grpo}
\emph{is} the mean field: each rollout is scored not on an absolute scale but against the \emph{group mean
baseline} $\operatorname{mean}(r)$ --- the Curie--Weiss $\avg{\cdot}$ --- exactly as each spin is scored against
the average magnetisation. The group is the ensemble.

\begin{center}
\begin{tabular}{@{}ll@{}}
\toprule
\textbf{Curie--Weiss variable} & \textbf{GRPO realisation} \\
\midrule
ensemble of $N$ spins & the group of $G$ rollouts for one prompt \\
mean field $\avg{\sigma}$ & the group-mean baseline $\operatorname{mean}(r)$ in \eqref{eq:grpo} \\
order parameter $m$ & the group's behavioural alignment (concentration of the $G$ rollouts) \\
external field $h$ & the reward $r$ \\
temperature $T$ & sampling temperature $\tau$ $+$ the $k3$ KL coefficient to the frozen reference \\
\bottomrule
\end{tabular}
\end{center}

The temperature has two contributions, both raising disorder. The sampling temperature $\tau$ sets the entropy
of the rollout distribution directly. The $k3$ KL coefficient to the frozen reference resists collapse toward a
single high-reward spike --- it pulls the group back toward the broad reference law, which is precisely a
temperature-like term keeping the ensemble from quenching. The reward is the field: it is what \emph{aims} the
update (Proposition~\ref{prop:aim}), and driving the advantage \eqref{eq:grpo} to a spike is the training-time
image of $m\to1$, the ordered phase. The group size $G$ plays the role of $N$: a larger group is a sharper
estimate of the mean field, i.e. a less finite-size-rounded transition.

\paragraph{The payoff: criticality as a hyperparameter schedule.} Because the ordered phase is over-committed
and the disordered phase learns nothing (Section~\ref{sec:control}), the GRPO knobs should place the group
\emph{near $T_c$}: diverse enough to explore (fluctuations survive, so a single high-reward rollout can still
move the group --- the training-time reading of surviving susceptibility) and coupled enough to converge (the
group-mean baseline is a meaningful signal, not noise). Concretely: choose a group size $G$ large enough that
$\operatorname{mean}(r)$ is a low-variance baseline; choose the sampling temperature $\tau$ and the KL
coefficient together so the group sits just below its critical point --- and then \emph{anneal} $T$ over
training, high early (explore) and cooling late (exploit), never quenching the advantage to a spike. A physics
analogy has become a schedule for three numbers the trainer already exposes. This is the same anneal the outer
loop prescribes for the population as a whole \cite{hanzo-evolutionary-orchestration}, read here at the level of
one GRPO group.

\section{Grounding II: the harness, the triple-duty critic, and the flywheel}
\label{sec:harness}

Inference-time, the job is to hold the running fleet at criticality. The single mechanism that makes this
practical is the adversarial critic, and its economy is that it does three jobs at once.

\begin{itemize}[leftmargin=1.4em,itemsep=2pt]
\item \textbf{Ferromagnetic fluctuation (anti-groupthink).} A critic whose standing instruction is ``try to
refute the current consensus'' is a source of thermal disorder aimed exactly where order is dangerous. Near
criticality (Proposition~\ref{prop:chi}) its large susceptibility means one well-aimed refutation can flip a bad
basin; it is the mechanism that keeps $T$ off the $m\to1$ floor.
\item \textbf{Reinforcement reward.} The same critic's verdict, as a scalar, is the reward $r$ of
Section~\ref{sec:grpo} --- the field that aims the policy update.
\item \textbf{Selection signal.} That verdict is also the gate that decides which trajectories survive the
round --- the selection operator of the outer evolutionary loop \cite{hanzo-evolutionary-orchestration}.
\end{itemize}

Crucially the critic's signal is largely \emph{ground truth}, not a learned reward model that can be gamed: a
green build compiles or it does not, an adversarial review passes or it does not, and the terminal signal is
whether the change ships and survives in production. Ground-truth reward is what lets the same scalar serve as
fluctuation, reward, and selection without any of the three being fooled by reward hacking. Winning
trajectories then distil back into \code{enso} by expert iteration \cite{hanzo-native-training}, folded in
through the continual-learning path that keeps task data private \cite{hanzo-continuous-learning-privacy}: a
better \code{enso} runs a better fleet, and a better fleet generates better trajectories.
\textbf{Model$\leftrightarrow$fleet co-evolution} is the flywheel, and it closes because one mechanism --- the
critic --- is simultaneously the thermostat's fluctuation, the trainer's reward, and the search's selection.

\paragraph{The three-timescale stack.} This paper builds the middle layer of a controller whose companion
\cite{hanzo-evolutionary-orchestration} builds the outer. Stacked, and separated by timescale so they compose
without fighting:
\begin{center}
\begin{tabular}{@{}lll@{}}
\toprule
\textbf{Layer} & \textbf{Timescale} & \textbf{Job} \\
\midrule
MFG inner (this paper, \S\ref{sec:iso}) & fast & each agent best-responds to the mean; the fleet equilibrates \\
Ferromagnet thermostat (this paper, \S\ref{sec:control}) & medium & hold $m$ near $T_c$; aim with the field $h$ \\
GA outer \cite{hanzo-evolutionary-orchestration} & slow & selection/recombination/mutation moves the population uphill \\
\bottomrule
\end{tabular}
\end{center}
One control variable --- $m$, the alignment / diversity of the fleet --- threads all three, and the whole craft
is annealing it from hot-diverse toward critical-aligned without quenching to $m=1$, aiming the descent with the
directive.

\section{What is real, what needs building, what is out of scope}
\label{sec:status}

\paragraph{\textsc{Real}.} The GRPO mean field is live and tested: the group-relative advantage
\eqref{eq:grpo}, its unbiased group-mean baseline, and the $k3$ KL term ship in
\code{ml/hanzo-training/src/grpo/} \cite{hanzo-native-training} with passing math and smoke tests. Read as
Section~\ref{sec:grpo} prescribes, GRPO already runs a Curie--Weiss ensemble --- the mean-field baseline is not
a metaphor bolted on, it is the code. The agentic harness already runs fleets of agents with adversarial critics
as gates, so the triple-duty critic of Section~\ref{sec:harness} is real in its selection and (via GRPO) reward
roles.

\paragraph{\textsc{Needs-Building}.} Two pieces make the thermostat explicit rather than implicit. First, an
\emph{$m$-estimator}: a cheap online measurement of the fleet's order parameter (behavioural concentration of
the agents' answers/trajectories) so the controller knows where on the $T$--$T_c$ axis it is sitting. Second, a
\emph{critic-gain controller} that raises critic strength as $m$ rises --- turning up adversarial pressure and
independent-lens diversity precisely when the fleet begins to over-order, and easing it as the fleet
fragments --- a closed loop that holds the operating point just below $T_c$. Both are orchestration logic over
the existing harness and trainer, not model or kernel work.

\paragraph{Out of scope.} This is a control lens, not a literal Hamiltonian. An agent swarm has no true energy
function; the physics buys the right \emph{variables} (order parameter, coupling, field, temperature) and the
phase-transition/criticality \emph{intuition}, and it earns the four theorems above about the reduced scalar
model --- but it does not hand us free rigor about the full swarm dynamics. We do not claim the fleet's
microstate is Gibbs-distributed, nor that $T_c=J$ holds numerically for real agents (the mean-field exponents
are the all-to-all idealisation); we claim the qualitative control law --- two phases, a critical operating
point, coherence-from-coupling and aim-from-field --- transfers, and that the GRPO grounding makes one instance
of it quantitative. Non-mean-field interaction structure (agents that couple locally rather than all-to-all),
finite-size corrections, and dynamical (out-of-equilibrium) fleet behaviour are left to the outer-loop treatment
\cite{hanzo-evolutionary-orchestration} and to future work.

\section{Conclusion}
\label{sec:conclusion}

A fleet run at zero coupling fragments; a fleet run at high coupling stampedes into a consensus it cannot
correct. The state we want --- coherent but correctable --- is not a fragile midpoint but a critical point, and
naming it correctly hands us the whole control law. A mean-field game closes best-response against
population-flow at a self-consistent equilibrium; Curie--Weiss ferromagnetism is that same equilibrium with a
scalar order parameter (Proposition~\ref{prop:isomorphism}), ordering below $T_c=J$
(Proposition~\ref{prop:tc}) and announcing the transition through a diverging susceptibility
(Proposition~\ref{prop:chi}). From those follow four knobs, one operating point (just below $T_c$), and one
sharp rule: coupling buys coherence but only the field aims it (Proposition~\ref{prop:aim}). The law is already
running in the stack --- GRPO's group-relative advantage \emph{is} the mean field, which turns criticality into
a hyperparameter schedule, and the harness's adversarial critic is at once the fleet's fluctuation, its reward,
and its selection, closing a model$\leftrightarrow$fleet co-evolution flywheel. What remains is to make the
thermostat explicit: measure $m$, and let critic gain rise with it, so the fleet holds itself at the edge of
chaos while the directive aims it at the objective.

\begin{thebibliography}{9}

\bibitem{lasrylions2007}
J.-M.\ Lasry and P.-L.\ Lions.
\newblock Mean field games.
\newblock \emph{Japanese Journal of Mathematics}, 2(1):229--260, 2007. (The HJB $\otimes$ Fokker--Planck
equilibrium.)

\bibitem{baxter1982}
R.\ J.\ Baxter.
\newblock \emph{Exactly Solved Models in Statistical Mechanics}.
\newblock Academic Press, 1982. (Ising / Curie--Weiss mean-field ferromagnetism and the self-consistent
magnetisation.)

\bibitem{shao2024grpo}
Z.\ Shao, P.\ Wang, Q.\ Zhu, R.\ Xu, J.\ Song, et al.
\newblock DeepSeekMath: Pushing the Limits of Mathematical Reasoning in Open Language Models.
\newblock \emph{arXiv:2402.03300}, 2024. (Group Relative Policy Optimization, GRPO.)

\bibitem{hanzo-native-training}
Hanzo AI --- ML Systems.
\newblock \emph{Native Training in hanzo-ml/hanzo-engine: One Engine, One Quant Format, One Backend}.
\newblock Hanzo Papers, 2026. (GRPO group-relative advantage, \code{ml/hanzo-training/src/grpo/}; expert
iteration.)

\bibitem{hanzo-evolutionary-orchestration}
Hanzo AI --- ML Systems.
\newblock \emph{Evolutionary Orchestration of Agentic Fleets: Population Search, Recombination, and the
Explore--Exploit Anneal}.
\newblock hanzo-agentic Paper 02, \emph{Continually Improving Fleets} series, 2026. Companion to this paper (the
outer evolutionary loop).

\bibitem{hanzo-consensus-zap-compute}
Hanzo AI --- ML Systems.
\newblock \emph{Consensus and ZAP Compute: A Distributed Substrate for Agentic Fleets}.
\newblock Hanzo Papers, 2026. (The shared-state / running-consensus substrate that carries the coupling.)

\bibitem{hanzo-continuous-learning-privacy}
Hanzo AI --- ML Systems.
\newblock \emph{Continuous Learning with Privacy}.
\newblock Hanzo Papers, 2026. (Folding winning trajectories back into the model while keeping task data
private.)

\end{thebibliography}

\end{document}
